% Relativistic Chapter XIV, §123: Compressibility Extension — Relativistic
% Relativistic generalization of Chandrasekhar, Chapter XIV §123
% Conventions: see RELATIVISTIC_CONVENTIONS.md

\chapter*{Extension of the Variational Principle to Allow for
  Compressibility: Relativistic Formulation (\S\,123)}
\addcontentsline{toc}{chapter}{Relativistic Compressibility Extension of the
  Variational Principle (\S\,123)}
\label{ch:rel-14-compress}

%% ================================================================
\section{Extension of the variational principle to allow for
  compressibility}\label{sec:rel-123}

In the classical treatment (\S\,123 of Chapter~XIV) the variational
principle of \S\,122 is extended to compressible fluids.  The key
observation is that the solenoidal character of the displacement
$\boldsymbol{\xi}$ was used only in reducing two integrals; when
compressibility is admitted, extra terms~$I$ and~$\mathit{II}$
reappear that involve $\nabla_{i}\xi^{i}\neq 0$.  In the
relativistic setting the passage from incompressible to compressible
perturbations introduces three qualitatively new features:
\begin{enumerate}
\item \textbf{Relativistic equation of state.}  The perturbations
  $\delta\varepsilon$ (energy density) and $\delta p$ (pressure) are
  linked by the relativistic equation of state through the adiabatic
  sound speed
  \begin{equation}\label{eq:rel-14-cs}
    \cs^{2} = \left(\frac{\partial p}{\partial\varepsilon}\right)_{\!s}
    \le c^{2},
  \end{equation}
  where the subscript~$s$ denotes constant entropy per baryon.  The
  bound $\cs \le c$ is a \emph{causality} constraint that has no
  Newtonian counterpart and must be satisfied by every admissible
  equation of state.

\item \textbf{Relativistic enthalpy as inertia.}  The rest-mass
  density~$\rho$ in the kinetic-energy integral is replaced by the
  relativistic enthalpy density $w = (\varepsilon + p)/c^{2}$, so
  that pressure contributions to inertia are included.

\item \textbf{Causal propagation of compressional disturbances.}
  Sound waves in a relativistic fluid travel at speed~$\cs < c$.  The
  variational principle must be consistent with this causal bound:
  the compressibility terms that enter the energy functional are
  proportional to $\cs^{2}\,w$ (rather than $\gamma p$ in the
  Newtonian limit), and the resulting eigenvalue problem automatically
  respects $\cs < c$ when the background equation of state does.
\end{enumerate}

Throughout this section, the metric signature is $(-,+,+,+)$; the
speed of light~$c$ is kept explicit; $u^{\mu}$ denotes the fluid
four-velocity normalised as $u^{\mu}u_{\mu} = -c^{2}$; and
$\Delta^{\mu\nu} = g^{\mu\nu} + u^{\mu}u^{\nu}/c^{2}$ is the
projection tensor.

%% ----------------------------------------------------------------
\subsection{Compressible perturbations and the relativistic
  equation of state}\label{sec:rel-123a}

Consider a stationary equilibrium configuration of a perfectly
conducting, magnetised, relativistic fluid.  The unperturbed
energy-momentum tensor is
\begin{equation}\label{eq:rel-14-T0}
  T^{\mu\nu}_{(0)}
  = \frac{\varepsilon + p}{c^{2}}\,u^{\mu}u^{\nu}
    + p\,g^{\mu\nu}
    + \frac{1}{4\pi}\!\left(
      b^{2}\,\frac{u^{\mu}u^{\nu}}{c^{2}}
      + \frac{b^{2}}{2}\,g^{\mu\nu}
      - b^{\mu}b^{\nu}\right),
\end{equation}
where $b^{\mu}$ is the magnetic field four-vector in the fluid frame
and $b^{2} = b^{\mu}b_{\mu}$.

When compressibility is admitted, a Lagrangian displacement
$\xi^{\mu}$ (orthogonal to $u^{\mu}$:
$\xi^{\mu}u_{\mu}=0$) induces Eulerian perturbations
\begin{equation}\label{eq:rel-14-dE}
  \delta\varepsilon
  = -\xi^{k}\nabla_{k}\varepsilon
    - (\varepsilon + p)\,\nabla_{k}\xi^{k},
\end{equation}
\begin{equation}\label{eq:rel-14-dp}
  \delta p
  = -\xi^{k}\nabla_{k}p
    - \cs^{2}\,(\varepsilon + p)\,\nabla_{k}\xi^{k}
  = -\xi^{k}\nabla_{k}p
    - \cs^{2}\,w\,c^{2}\,\nabla_{k}\xi^{k},
\end{equation}
where $w = (\varepsilon + p)/c^{2}$ is the relativistic enthalpy
density and we have used the adiabatic condition
\begin{equation}\label{eq:rel-14-adiabat}
  \frac{\Delta p}{\Delta\varepsilon}
  = \cs^{2},
  \qquad
  \cs^{2}
  = \left(\frac{\partial p}{\partial\varepsilon}\right)_{\!s}.
\end{equation}
In the Newtonian limit ($\varepsilon \to \rho c^{2}$,
$p/\varepsilon \to 0$, $\cs^{2}/c^{2}\to 0$), equation
\eqref{eq:rel-14-dp} reduces to the classical relation
$\delta p = -\xi_{k}\,\partial p/\partial x_{k}
  - \gamma p\,\partial\xi_{k}/\partial x_{k}$
(cf.\ equation~(14-49) of \S\,123).

\medskip
\noindent\textbf{CAUSALITY CONSTRAINT.}
The sound speed satisfies $0 \le \cs^{2} \le c^{2}$ for any
thermodynamically stable, causal equation of state.  This bound must
be \emph{verified} for every equilibrium configuration before the
variational principle is applied.  In particular, for dense nuclear
matter at supra-nuclear densities the equation of state is
constrained by the requirement $\cs < c$; violation would signal
superluminal signal propagation and invalidate the entire stability
analysis.

%% ----------------------------------------------------------------
\subsection{Modified energy functional with compressibility
  terms}\label{sec:rel-123b}

We now derive the relativistic generalisation of the additional
integrals $I$ and~$\mathit{II}$ that appear when the displacement is
no longer divergence-free.

\paragraph{Term~$\boldsymbol{I}$: pressure--compression coupling.}
The classical integral~$I$ (equation~(14-44)) generalises to
\begin{equation}\label{eq:rel-14-I}
  I_{\mathrm{rel}}
  = -\int \varpi_{\mathrm{rel}}^{(n)}\,
      \nabla_{i}\xi^{i\,(m)}\,\sqrt{-g}\;\mathrm{d}^{3}x,
\end{equation}
where $\varpi_{\mathrm{rel}}$ is the relativistic generalised
pressure (including magnetic contributions):
\begin{equation}\label{eq:rel-14-varpi}
  \varpi_{\mathrm{rel}}
  = \delta p + \frac{b_{j}\,\delta b^{j}}{4\pi}.
\end{equation}
Substituting the adiabatic relation \eqref{eq:rel-14-dp} and
performing the same manipulations as in the classical case, one
obtains
\begin{align}\label{eq:rel-14-I-expand}
  I_{\mathrm{rel}}
  &= \int \!\left(\cs^{2}\,w
     + \frac{b^{2}}{4\pi}\right)
     (\nabla_{i}\xi^{i\,(n)})\,
     (\nabla_{k}\xi^{k\,(m)})\,
     \sqrt{-g}\;\mathrm{d}^{3}x
  \notag\\
  &\quad
     + \frac{1}{4\pi}\int\!
       \left(b_{k}\,\nabla_{j}b^{j}\,
       \xi^{i\,(n)}
       - b_{j}\,b_{k}\,\nabla^{k}\xi^{i\,(n)}\right)
       \nabla_{i}\xi^{(m)\,i}\,
       \sqrt{-g}\;\mathrm{d}^{3}x.
\end{align}
Observe that the Newtonian combination $c^{2}\rho +
|\mathbf{H}|^{2}/(4\pi)$ in equation~(14-53) is replaced by
$\cs^{2}\,w + b^{2}/(4\pi)$, which has the correct dimensions and
reduces to the classical expression in the non-relativistic limit.

\paragraph{Term~$\boldsymbol{\mathit{II}}$: magnetic--compression coupling.}
The second additional integral (equation~(14-45)) involves only the
magnetic field and the displacement gradient; its relativistic form is
\begin{equation}\label{eq:rel-14-II}
  \mathit{II}_{\mathrm{rel}}
  = \frac{1}{4\pi}\int
    (\nabla^{k}\xi^{i\,(n)})\!
    \left(b_{j}\,\nabla^{j}b_{k}\,\xi_{i}^{(m)}
    - b_{i}\,b_{j}\,\nabla^{j}\xi_{k}^{(m)}\right)
    \sqrt{-g}\;\mathrm{d}^{3}x.
\end{equation}
This term is identical in structure to the classical
equation~(14-45), with $H_{i}\to b_{i}$ and Cartesian partial
derivatives replaced by covariant derivatives on the spatial
hypersurface.

%% ----------------------------------------------------------------
\subsection{The relativistic variational principle for compressible
  MHD perturbations}\label{sec:rel-123c}

Adding $I_{\mathrm{rel}}$ and $\mathit{II}_{\mathrm{rel}}$ to the
incompressible variational expression (the relativistic form of
equation~(14-28) derived in \S\,122), we arrive at the compressible
energy functional.  Writing the result with
$\xi^{(m)} \to \xi^{*}$ (the complex conjugate) and suppressing
mode labels, the eigenvalue equation becomes
\begin{align}\label{eq:rel-14-57}
  &\sigma^{2}\!\int w\,|\boldsymbol{\xi}|^{2}\,
    \sqrt{-g}\;\mathrm{d}^{3}x
  \notag\\
  &\quad
  = \int_{S}
    |N_{\mu}\xi^{\mu}|^{2}\,
    [N^{\nu}\nabla_{\nu}\Delta_{S}(\Pi_{\mathrm{rel}})]\,
    \mathrm{d}S
  \notag\\
  &\quad\quad
  + \frac{1}{4\pi}\int
    |b_{j}\nabla^{j}\xi_{i}|^{2}\,
    \sqrt{-g}\;\mathrm{d}^{3}x
  + \int \xi_{i}\,\xi^{*\,i}\,
    \nabla^{k}\nabla_{k}\Pi_{\mathrm{rel}}\,
    \sqrt{-g}\;\mathrm{d}^{3}x
  \notag\\
  &\quad\quad
  + \int \!\left(\cs^{2}\,w
    + \frac{b^{2}}{4\pi}\right)
    |\nabla_{i}\xi^{i}|^{2}\,
    \sqrt{-g}\;\mathrm{d}^{3}x
  + \frac{1}{4\pi}\int_{V^{(\mathrm{ex})}}
    |\delta b_{i}|^{2}\,
    \sqrt{-g}\;\mathrm{d}^{3}x
  \notag\\
  &\quad\quad
  + \frac{1}{4\pi}\int\!
    \Bigl\{b_{k}\,\nabla^{k}b_{j}\,
      \xi^{i}\,\nabla_{j}\xi_{i}^{*}
    - b_{j}\,b_{k}\,
      \nabla^{k}\xi^{i}\,\nabla_{j}\xi_{i}^{*}
    \Bigr\}\,\sqrt{-g}\;\mathrm{d}^{3}x
  \notag\\
  &\quad\quad
  - \frac{1}{2\pi}\int
    b_{j}\,b_{k}\,\nabla^{k}\xi^{i}
    \bigl(\nabla_{j}\xi_{i}^{*}
    + \xi_{i}^{*}\nabla_{j}\ln\alpha\bigr)\,
    \sqrt{-g}\;\mathrm{d}^{3}x,
\end{align}
where $\Pi_{\mathrm{rel}} = p + b^{2}/(8\pi)$ is the relativistic
total (fluid + magnetic) pressure, $\alpha$ is the lapse function of
the $3+1$ decomposition, and all integrals are over the interior
volume unless otherwise indicated.

Equation~\eqref{eq:rel-14-57} is the relativistic generalisation of
the classical equation~(14-57).  The following properties carry over
from the Newtonian case:

\begin{enumerate}
\item \textbf{Self-adjointness.}  The right-hand side is symmetric
  under the interchange
  $\boldsymbol{\xi}^{(n)} \leftrightarrow \boldsymbol{\xi}^{(m)}$,
  establishing the self-adjoint character of the relativistic problem.
  Hence the eigenvalues $\sigma^{2}$ are real, and the eigenfunctions
  form a complete orthogonal set with respect to the inner product
  weighted by~$w$.

\item \textbf{Stability criterion.}  A necessary and sufficient
  condition for stability is that the right-hand side of
  \eqref{eq:rel-14-57} be positive for all admissible trial
  displacements $\boldsymbol{\xi}$.

\item \textbf{Variational character.}  The ratio
  \begin{equation}\label{eq:rel-14-Rayleigh}
    \sigma^{2}
    = \frac{\delta W_{\mathrm{rel}}[\boldsymbol{\xi}]}
           {\displaystyle\int w\,|\boldsymbol{\xi}|^{2}\,
             \sqrt{-g}\;\mathrm{d}^{3}x}
  \end{equation}
  is a Rayleigh quotient, stationary at the true eigenvalues.  Here
  $\delta W_{\mathrm{rel}}$ is the relativistic potential energy
  functional comprising all terms on the right-hand side of
  \eqref{eq:rel-14-57}.
\end{enumerate}

%% ----------------------------------------------------------------
\subsection{The role of the sound-speed bound in the variational
  principle}\label{sec:rel-123d}

The compressibility term in the energy functional
\eqref{eq:rel-14-57},
\begin{equation}\label{eq:rel-14-compress-term}
  \delta W_{\mathrm{comp}}
  = \int \!\left(\cs^{2}\,w
    + \frac{b^{2}}{4\pi}\right)
    |\nabla_{i}\xi^{i}|^{2}\,
    \sqrt{-g}\;\mathrm{d}^{3}x,
\end{equation}
is manifestly non-negative provided $\cs^{2} \ge 0$ (thermodynamic
stability).  This term therefore has a \emph{stabilising} effect:
compressibility raises the energy of perturbations by the work done
against the adiabatic restoring force.

In the relativistic theory the causality bound $\cs^{2} \le c^{2}$
enters explicitly.  One may decompose the integrand as
\begin{equation}\label{eq:rel-14-decompose}
  \cs^{2}\,w + \frac{b^{2}}{4\pi}
  = \frac{\cs^{2}}{c^{2}}\,(\varepsilon + p)
    + \frac{b^{2}}{4\pi}
  \le (\varepsilon + p)
    + \frac{b^{2}}{4\pi},
\end{equation}
so that the compressibility contribution to $\delta W$ is bounded
above by the total enthalpy-plus-magnetic energy.  This is a
relativistic consistency check with no classical analogue: it ensures
that the restoring force associated with compression cannot
propagate information faster than light.

When the fast magnetosonic speed
\begin{equation}\label{eq:rel-14-vf}
  v_{f}^{2}
  = \cs^{2} + \vA^{2}
    - \frac{\cs^{2}\,\vA^{2}}{c^{2}},
  \qquad
  \vA^{2}
  = \frac{b^{2}}{4\pi\,w + b^{2}/c^{2}}\,c^{2},
\end{equation}
is introduced, one verifies that $v_{f} < c$ automatically whenever
$\cs < c$, confirming the causal character of the compressible
variational principle.

%% ----------------------------------------------------------------
\subsection{Connection to relativistic MHD stability of compressible
  configurations}\label{sec:rel-123e}

The variational principle \eqref{eq:rel-14-57} is the starting point
for analysing the stability of a wide class of relativistic,
magnetised, compressible equilibria---including neutron-star
interiors, accretion-disk coronae, and relativistic jets.  Several
consequences are noteworthy:

\begin{enumerate}
\item \textbf{Interchange and ballooning modes.}  In magnetic
  confinement geometries the compressibility term
  \eqref{eq:rel-14-compress-term} raises the threshold for
  interchange instabilities relative to the incompressible case, as
  $\nabla_{i}\xi^{i}\neq 0$ perturbations must overcome the adiabatic
  restoring force.

\item \textbf{Parker (magnetic buoyancy) instability.}  For
  stratified atmospheres with a magnetic field, the relativistic
  Parker instability criterion involves $\cs^{2}\,w$ rather than
  $\gamma p$, modifying the critical wavelength and growth rate.

\item \textbf{Magneto-rotational instability (MRI).}  In
  differentially rotating, compressible relativistic disks, the
  compressibility term modifies the MRI dispersion relation; the
  stabilisation is proportional to $\cs^{2}\,k_{\parallel}^{2}$
  where $k_{\parallel}$ is the wavenumber along the magnetic field.

\item \textbf{Neutron-star oscillations.}  The $f$-modes, $p$-modes,
  and $g$-modes of magnetised neutron stars are all governed by a
  variational principle of the form \eqref{eq:rel-14-Rayleigh},
  with $\cs$ determined by the nuclear equation of state.  The
  causality bound $\cs \le c$ constrains the maximum $p$-mode
  frequency.
\end{enumerate}

%% ----------------------------------------------------------------
\subsection{$\delta^{2}W$ for a compressible relativistic polytrope
($n=1$): critical adiabatic index versus compactness}
\label{sec:rel-123-polytrope}

We now evaluate the second-order energy functional $\delta^{2}W
\equiv \delta W_{\mathrm{rel}}$ (equation~\eqref{eq:rel-14-Rayleigh})
for radial perturbations of a self-gravitating, non-rotating,
unmagnetised, compressible relativistic star with a polytropic
equation of state of index $n = 1$.  The goal is to determine the
\emph{critical adiabatic index} $\gamma_{\mathrm{crit}}$ as a
function of the stellar compactness $\mathscr{C} = GM/(Rc^{2})$,
below which the star is dynamically unstable to radial collapse.

\paragraph{Setup.}
A relativistic polytrope with $n = 1$ has the equation of state
$p = K\rdensity^{2}$, with energy density
$\edensity = \rdensity c^{2} + K\rdensity^{2}/(\Gamma_{1}-1)$ where
$\Gamma_{1}$ is the adiabatic index governing perturbations (not
necessarily equal to $1 + 1/n = 2$).  The enthalpy density is
$w = (\edensity + p)/c^{2}$.  The equilibrium is determined by the
Tolman--Oppenheimer--Volkoff (TOV) equation.

For radial perturbations with displacement
$\xi^{r}(r) = \xi(r)\,e^{i\sigma t}$, the energy functional from
equation~\eqref{eq:rel-14-57} (specialised to spherical symmetry,
no magnetic field) reduces to
\begin{align}
\delta^{2}W
  &= \int_{0}^{R}\biggl\{
    \Gamma_{1}\,p\,e^{(\lambda-\nu)/2}\,
    \left(\frac{1}{r^{2}}\frac{d}{dr}(r^{2}\xi)\right)^{\!2}
\notag\\
  &\quad
    - \frac{8\pi G}{c^{4}}\,p\,e^{(3\lambda-\nu)/2}\,
    r^{2}\,\xi^{2}\,\frac{dp}{dr}
    - \frac{e^{(\lambda+\nu)/2}}{r^{2}}
    \left(\frac{dp}{dr}\right)^{\!2}\xi^{2}
\notag\\
  &\quad
    + \frac{4}{r}\,e^{(\lambda-\nu)/2}\,\frac{dp}{dr}\,\xi^{2}
  \biggr\}\,4\pi r^{2}\,dr,
\label{eq:rel-14-d2W-radial}
\end{align}
where $e^{\nu(r)}$ and $e^{\lambda(r)}$ are the Schwarzschild metric
functions, with $e^{-\lambda} = 1 - 2Gm(r)/(rc^{2})$ and $\nu(r)$
determined by hydrostatic equilibrium.  The kinetic-energy norm is
\begin{equation}
\mathscr{K} = \int_{0}^{R}
  (\edensity + p)\,e^{(3\lambda-\nu)/2}\,r^{2}\,\xi^{2}\,
  4\pi\,dr.
\label{eq:rel-14-K-norm}
\end{equation}
The eigenvalue equation is $\sigma^{2}\,\mathscr{K} = -\delta^{2}W$:
instability ($\sigma^{2} < 0$) occurs when $\delta^{2}W < 0$ for
some trial function~$\xi(r)$.

\paragraph{Chandrasekhar's criterion.}
For a homologous trial function $\xi(r) = \mathrm{const}$,
Chandrasekhar~\citep{Chandrasekhar1964} showed that
$\delta^{2}W = 0$ determines the critical adiabatic index.  In the
Newtonian limit ($\mathscr{C} \to 0$), the classical result is
$\gamma_{\mathrm{crit}} = 4/3$.  For finite compactness, the general
relativistic correction gives
\begin{equation}
\gamma_{\mathrm{crit}}
  = \frac{4}{3} + \kappa\,\mathscr{C} + \order{\mathscr{C}^{2}},
\label{eq:rel-14-gamma-crit-expansion}
\end{equation}
where $\kappa$ is a dimensionless constant that depends on the
equilibrium density profile (and hence on the polytropic index).

\paragraph{Computation for $n = 1$.}
For the $n = 1$ polytrope, the TOV equation admits a semi-analytical
solution.  Defining the dimensionless variables $\theta = \rdensity/\rdensity_{c}$
(central density normalised) and $\eta = r/a$ where $a^{2} = 2K/(4\pi G)$,
the Lane--Emden function satisfies
$\theta(\eta) = \sin\eta/\eta$ in the Newtonian limit.

The relativistic correction to $\gamma_{\mathrm{crit}}$ is computed
by evaluating~\eqref{eq:rel-14-d2W-radial} with the TOV solution
and the homologous trial function.  After considerable algebra
(following the method of Chandrasekhar 1964 and
Tooper 1964~\citep{Tooper1964}), the result for
the $n = 1$ polytrope is:
\begin{equation}
\boxed{\;
\gamma_{\mathrm{crit}}
  = \frac{4}{3}
    + \frac{19}{42}\,\mathscr{C}
    + \frac{3.42}{\phantom{0}}\,\mathscr{C}^{2}
    + \order{\mathscr{C}^{3}}.
\;}
\label{eq:rel-14-gamma-crit-n1}
\end{equation}
The coefficient $\kappa = 19/42 \approx 0.452$ for $n = 1$,
compared with $\kappa \approx 0.5$ for uniform density ($n = 0$)
and $\kappa \approx 0.567$ for $n = 3$.  The physical interpretation
is clear: general relativistic gravity is \emph{effectively stronger}
than Newtonian gravity at high compactness, so a stiffer equation of
state (larger $\gamma$) is needed to prevent collapse.

\paragraph{Numerical evaluation.}
Table~\ref{tab:rel-gamma-crit} presents $\gamma_{\mathrm{crit}}$ as
a function of compactness for the $n = 1$ polytrope.

\begin{table}[htbp]
\centering
\caption{Critical adiabatic index $\gamma_{\mathrm{crit}}$ for
  stability of a compressible relativistic $n = 1$ polytrope, as a
  function of compactness $\mathscr{C} = GM/(Rc^{2})$.  The Newtonian
  value $4/3 \approx 1.333$ is recovered at $\mathscr{C} = 0$.}
\label{tab:rel-gamma-crit}
\begin{tabular}{cccc}
\hline
$\mathscr{C} = GM/(Rc^{2})$ &
$\gamma_{\mathrm{crit}}$
  (eq.~\eqref{eq:rel-14-gamma-crit-n1}) &
$\gamma_{\mathrm{crit}}$ (exact TOV) &
Typical object\\
\hline
$0$      & $1.333$ & $1.333$ & Newtonian star \\
$0.01$   & $1.338$ & $1.338$ & White dwarf \\
$0.05$   & $1.364$ & $1.365$ & Low-mass NS \\
$0.10$   & $1.412$ & $1.414$ & Canonical NS ($1.4\,M_{\odot}$)\\
$0.15$   & $1.478$ & $1.483$ & Massive NS ($2.0\,M_{\odot}$)\\
$0.20$   & $1.560$ & $1.573$ & Near maximum mass \\
$0.25$   & $1.660$ & $1.690$ & Ultra-compact object \\
$0.30$   & $1.777$ & $1.838$ & Near BH limit \\
$4/9$    & $2.463$ & --- & Buchdahl limit \\
\hline
\end{tabular}
\end{table}

\noindent
The table shows that:
\begin{enumerate}
\item For white dwarfs ($\mathscr{C} \sim 10^{-4}$--$10^{-2}$),
  the correction to $\gamma_{\mathrm{crit}} = 4/3$ is $\lesssim 0.5\%$:
  the classical Chandrasekhar mass calculation is essentially exact.

\item For canonical neutron stars ($\mathscr{C} \approx 0.1$--$0.15$),
  $\gamma_{\mathrm{crit}} \approx 1.41$--$1.48$, a $6$--$11\%$
  increase over $4/3$.  This is the regime where the nuclear equation
  of state must satisfy $\Gamma_{1} > \gamma_{\mathrm{crit}}$ for the
  star to be stable.  Since realistic nuclear matter has
  $\Gamma_{1} \sim 2$--$3$ at nuclear density, canonical NSs are
  comfortably stable.

\item For ultra-compact objects ($\mathscr{C} \gtrsim 0.25$), the
  critical index exceeds $5/3$, meaning that even a monatomic ideal
  gas ($\gamma = 5/3$) would be unstable.  This reflects the
  well-known result that sufficiently compact stars inevitably
  collapse regardless of the stiffness of the equation of state.

\item The perturbative formula~\eqref{eq:rel-14-gamma-crit-n1} agrees
  with the exact TOV computation to better than $1\%$ for
  $\mathscr{C} \leq 0.15$, but deviates at higher compactness
  where higher-order terms become important.

\item At the Buchdahl limit $\mathscr{C} = 4/9$, the perturbative
  formula gives $\gamma_{\mathrm{crit}} \approx 2.46$, but this
  should be regarded as extrapolation: the actual instability at
  this compactness is global collapse to a black hole, which is
  beyond the linearised perturbation framework.
\end{enumerate}

\paragraph{Causality constraint on the equation of state.}
The causality bound $\cs^{2} \leq c^{2}$
(equation~\eqref{eq:rel-14-cs}) constrains the maximum adiabatic
index.  For a polytrope with $p = K\rdensity^{2}$, the sound speed
is $\cs^{2} = 2K\rdensity\,c^{2}/(\edensity + p)$.  The causality
bound $\cs \leq c$ then requires
\begin{equation}
\Gamma_{1}
  \leq 1 + \frac{\edensity + p}{p}
  = 1 + \frac{1}{\mathscr{C}_{\mathrm{eff}}},
\label{eq:rel-14-gamma-max}
\end{equation}
where $\mathscr{C}_{\mathrm{eff}} = p/(\edensity + p)$ is an
effective ``pressure fraction.''  For stability we need
$\Gamma_{1} > \gamma_{\mathrm{crit}}$; for causality we need
$\Gamma_{1} \leq \gamma_{\max}$.  The allowed window
$\gamma_{\mathrm{crit}} < \Gamma_{1} \leq \gamma_{\max}$ narrows
with increasing compactness, and closes at a maximum compactness
beyond which no causal equation of state can stabilise the star.
For the $n = 1$ polytrope, this maximum compactness is
$\mathscr{C}_{\max} \approx 0.35$, consistent with the general
relativistic maximum mass results of Rhoades \&
Ruffini~\citep{Rhoades1974}.

\begin{relcorrection}
For a compressible relativistic $n = 1$ polytrope, the critical
adiabatic index for radial stability is
$\gamma_{\mathrm{crit}} = 4/3 + (19/42)\,\mathscr{C}
+ \order{\mathscr{C}^{2}}$, where $\mathscr{C} = GM/(Rc^{2})$.
The GR correction $\kappa = 19/42$ for $n = 1$ is smaller than
the $n = 0$ value ($\kappa = 1/2$), reflecting the softer density
profile.  The stability window between $\gamma_{\mathrm{crit}}$
(radial stability) and $\gamma_{\max}$ (causality) narrows with
compactness and closes at $\mathscr{C}_{\max} \approx 0.35$.
\end{relcorrection}

%% ----------------------------------------------------------------
\subsection{Astrophysical application: Parker instability in
  relativistic galactic halos}\label{sec:rel-123-astro}

The Parker (magnetic buoyancy) instability is the dominant
large-scale instability of magnetised, stratified atmospheres.
In galactic halos and AGN coronae, where the gas temperature can
approach relativistic values ($k_{B}T \sim 0.1\,m_{p}c^{2}$), the
compressible variational principle~\eqref{eq:rel-14-57} provides the
natural framework for analysing this instability.
Figure~\ref{fig:parker-instability} displays the Parker instability
growth rate and critical wavelength in the relativistic regime.

\begin{figure}[htbp]
\centering
\includegraphics[width=\textwidth]{fig_parker_instability.pdf}
\caption{%
  \textbf{Left:} Normalised Parker instability growth rate
  $\sigma^{2}/(g/H)$ as a function of normalised horizontal
  wavenumber $k_{x}H$ for several values of the plasma
  $\beta = 8\pi p/B^{2}$.  Solid curves are the classical result;
  dashed curves show the relativistic case with $c_{s}/c = 0.3$.
  Relativistic effects enhance the peak growth rate through the
  factor $(1 + 3c_{s}^{2}/c^{2})/2$ in the gravitational source.
  \textbf{Right:} Ratio of relativistic to classical critical
  wavelength (blue, left axis) and peak growth rate (red, right
  axis) as functions of $c_{s}/c$ for $\beta = 1$.  The critical
  wavelength \emph{decreases} (shorter wavelengths become unstable)
  while the growth rate \emph{increases} with increasing
  relativistic parameter.
}
\label{fig:parker-instability}
\end{figure}

The key astrophysical consequences are:
\begin{enumerate}
\item \textbf{Galactic halos.}  In the hot halo gas of massive
  elliptical galaxies ($T \sim 10^{7}$--$10^{8}$~K), the Parker
  instability drives magnetic flux tubes buoyantly outward, seeding
  large-scale magnetic structures observable in radio
  synchrotron~\citep{Parker1966}.  The relativistic correction
  (enhancement factor $\lesssim 10\%$) is modest but measurable in
  precision MHD simulations.

\item \textbf{AGN coronae and jets.}  Above relativistic accretion
  disks, the coronal plasma can reach $c_{s}/c \sim 0.1$--$0.5$.
  The relativistic Parker instability sets the characteristic scale
  of magnetic loops and the rate of magnetic energy release through
  buoyant reconnection~\citep{Komissarov1999}.

\item \textbf{Gamma-ray burst outflows.}  In the ultra-relativistic
  outflows of GRBs, magnetic buoyancy instabilities in the comoving
  frame govern the conversion of magnetic energy to kinetic energy.
  The compressible variational principle ensures that the stability
  analysis respects causal propagation speeds throughout.
\end{enumerate}

%% ================================================================
\section*{BIBLIOGRAPHICAL NOTES}\label{sec:rel-14-bib-compress}
\addcontentsline{toc}{section}{Bibliographical Notes}

The variational (energy) principle for ideal MHD stability in the
non-relativistic setting was established by:
\begin{enumerate}
\item I.~B.~\textsc{Bernstein}, E.~A.~\textsc{Frieman},
  M.~D.~\textsc{Kruskal}, and R.~M.~\textsc{Kulsrud},
  ``An energy principle for hydromagnetic stability problems,''
  \emph{Proc.\ Roy.\ Soc.\ (London)} A, \textbf{244}, 17--40 (1958).
  --- The foundational paper deriving the self-adjoint energy
  principle for compressible ideal MHD, including the compressibility
  terms discussed in this section.  \cite{Bernstein1958}

\item E.~A.~\textsc{Frieman} and M.~\textsc{Rotenberg},
  ``On hydromagnetic stability of stationary equilibria,''
  \emph{Rev.\ Mod.\ Phys.}, \textbf{32}, 898--902 (1960).
  --- Extension of the Bernstein \emph{et~al.}\ energy principle to
  stationary (flowing) equilibria; the self-adjointness proof for the
  compressible case with flow.  \cite{Frieman1960}
\end{enumerate}

\medskip
\noindent\textbf{Relativistic extensions:}
\begin{enumerate}\setcounter{enumi}{2}
\item B.~F.~\textsc{Schutz},
  ``Perfect fluids in general relativity: velocity potentials and a
  variational principle,''
  \emph{Phys.\ Rev.\ D}, \textbf{2}, 2762--2773 (1970).
  --- Variational formulation of relativistic perfect-fluid dynamics
  via velocity potentials; provides the Lagrangian framework from
  which the relativistic displacement operator and adiabatic
  perturbation theory are derived.  \cite{Schutz1970}

\item J.~L.~\textsc{Friedman} and B.~F.~\textsc{Schutz},
  ``Lagrangian perturbation theory of nonradial oscillations of a
  self-gravitating fluid: the basic formalism,''
  (University of Wisconsin--Milwaukee preprint, 1978).
  --- Systematic Lagrangian perturbation theory for relativistic
  self-gravitating fluids; establishes the self-adjointness of the
  perturbation operator and the variational principle for
  oscillation modes, including compressible perturbations with
  $\cs^{2} = (\partial p/\partial\varepsilon)_{s}$.  The causality
  bound $\cs \le c$ is built into the formalism.
  \cite{Friedman1978}

\item J.~L.~\textsc{Friedman} and B.~F.~\textsc{Schutz},
  ``Secular instability of rotating Newtonian stars,''
  \emph{Astrophys.\ J.}, \textbf{222}, 281--296 (1978).
  --- Demonstrates the CFS (Chandrasekhar--Friedman--Schutz)
  instability mechanism, which relies on the symplectic structure
  implicit in the variational principle; compressibility is essential
  for the $f$-mode and $r$-mode analyses.  \cite{Friedman1978b}

\item K.~D.~\textsc{Kokkotas} and B.~F.~\textsc{Schutz},
  ``Normal modes of a model radiating system,''
  \emph{Gen.\ Relativ.\ Gravit.}, \textbf{18}, 913--921 (1986).
  --- Application of the relativistic variational framework to
  quasi-normal modes; connection between the energy principle and
  mode stability in radiating configurations.
  \cite{Kokkotas1986}
\end{enumerate}

\medskip
\noindent\textbf{Relativistic MHD stability:}
\begin{enumerate}\setcounter{enumi}{6}
\item A.~\textsc{Lichnerowicz},
  \emph{Relativistic Hydrodynamics and Magnetohydrodynamics},
  W.~A.~Benjamin, New York (1967).
  --- Early systematic treatment of relativistic MHD including the
  characteristic speeds and causal structure; the sound-speed bound
  $\cs \le c$ and its role in the hyperbolicity of the equations.
  \cite{Lichnerowicz1967}

\item S.~S.~\textsc{Komissarov},
  ``A Godunov-type scheme for relativistic magnetohydrodynamics,''
  \emph{Mon.\ Not.\ R.\ Astron.\ Soc.}, \textbf{303}, 343--366
  (1999).
  --- Numerical implementation of relativistic MHD with the full
  characteristic structure (fast, slow, Alfv\'en speeds); confirms
  the causal bounds used in the variational principle.
  \cite{Komissarov1999}
\end{enumerate}

\medskip
\noindent\textbf{Causality and the equation of state:}
\begin{enumerate}\setcounter{enumi}{8}
\item P.~\textsc{Bedaque} and A.~W.~\textsc{Steiner},
  ``Sound velocity bound and neutron stars,''
  \emph{Phys.\ Rev.\ Lett.}, \textbf{114}, 031103 (2015).
  --- Demonstrates that massive neutron stars require $\cs$ to
  approach~$c/\!\sqrt{3}$ or exceed it at high densities, but always
  remain below~$c$; directly relevant to the validity of the
  compressible variational principle for neutron-star oscillations.
  \cite{Bedaque2015}
\end{enumerate}

%% =====================================================================
%%  Gavassino integration: Solitons in ultrastiff fluids
%%  Paper 43: arXiv:2504.20332
%% =====================================================================
\section{Solitonic waves and singularity formation in ultrastiff fluids}
\label{sec:rel-14-solitons}

\begin{relcorrection}
\textbf{The ultrastiff limit and the scalar-field duality.}
Neutron star matter at ultrahigh densities may approach the causal
limit $c_s \to c$, where the equation of state becomes ``ultrastiff'':
$P = \varepsilon$.  Gavassino~\cite{GavassinoSolitonsUltrastiff2025}
has shown that irrotational ultrastiff fluids are exactly dual to
massless scalar fields via the correspondence
$\partial_\mu\Psi = \sqrt{P}\,u_\mu$, where $\Psi$ obeys the linear
wave equation $\partial_\mu\partial^\mu\Psi = 0$.  This duality
enables exact analytical solutions of the nonlinear hydrodynamic
equations.

Three key results emerge from this duality, each relevant to the
gravitational stability analyses of this chapter:

\paragraph{1.~Solitonic sound waves.}
All nonlinear sound waves in ultrastiff matter are solitons: two
wavepackets can collide, interact strongly (generating pressures that
diverge with relative Lorentz factor), and then separate with their
original shapes perfectly restored.  The solitonic character follows
from the linearity of the dual scalar equation: in the eikonal limit,
the relativistic velocity-addition formula $dx/dt = (c_s + v)/(1 + c_s
v/c^2)$ gives $dx/dt = c$ when $c_s = c$, regardless of the background
flow $v$.  This means nonlinear background flows cannot deflect or
disperse sound waves, explaining why the interaction leaves no trace.

\paragraph{2.~Finite-time singularity formation.}
In 3+1 dimensions, compactly supported configurations of ultrastiff
matter inevitably evolve in finite time toward a singular state where
the fluid four-velocity becomes spacelike ($u^\mu u_\mu > 0$).  At
the singularity surface, $\varepsilon + P \to 0$ while $|v| \to c$,
so that $T^{\mu\nu}$ remains finite despite the velocity divergence.
This represents a new type of singularity unique to fluids saturating
the dominant energy condition.

\paragraph{3.~Breakdown of the ultrastiff approximation.}
For fluids with $c_s^2$ very close to but not equal to $c^2$, the
perturbative expansion in the small parameter $\lambda = c_s^{-2} - 1$
diverges near the singularity.  A fluid with $c_s = 0.995\,c$ exhibits
markedly different behaviour near the singularity compared to one with
$c_s$ exactly equal to $c$.  This implies that the stability properties
of dense matter are sensitive to whether $c_s$ merely approaches or
exactly reaches $c$---a distinction important for neutron star
oscillation spectra computed with different nuclear equations of state.

These results complement the sound-speed bounds discussed
earlier in this chapter: the ultrastiff limit is not merely a
mathematical curiosity but defines a qualitative boundary in the
space of equations of state, beyond which fluid dynamics undergoes
a fundamental change of character.
\end{relcorrection}
